\section{Preliminaries and Related Work}
\label{ch:related-work}

To provide the necessary background for the evaluation of the xLSTM-7B architecture,
this chapter presents the theoretical concepts and related work relevant to this study.
The first section reviews the evolution of recurrent neural networks, from classical RNNs
to the modern architecture proposed by Beck et al. \cite{beck2024xlstmextendedlongshortterm} 
that incorporate gating mechanisms and matrix-based memory structures. \\ \\
%
\subsection{Recurrent Neural Networks}
Recurrent Neural Networks (RNNs) represent one of the earliest neural architectures specifically 
designed for sequence modeling tasks. Their foundations can be traced to the work of Rumelhart et al. 
\cite{rumelhart1986learning}, who introduced the backpropagation algorithm and discussed neural networks 
with recurrent connections. Unlike feedforward neural networks, which process each input independently, 
RNNs maintain a hidden state that evolves over time, allowing information from previous inputs to influence 
future predictions. \\ \\
%
The key idea behind recurrent architectures is the reuse of the same set of parameters at every time step. 
Given an input sequence, the hidden state is updated recursively based on both the current input and the 
previous hidden state, enabling the network to capture temporal dependencies within the data. This property 
made RNNs a natural choice for tasks involving sequential information, such as language modeling, speech 
recognition, machine translation, and time-series forecasting.

\begin{figure}[htbp]
    \centering
    \includegraphics[width=\textwidth]{RNN.png}
    \caption{RNN: compressed (left) and unfolded(right), diagram inspired by \cite{olah2015lstm}}
    \label{fig:RNN-diagram}
\end{figure}
\noindent The operation of a standard RNN can be formally described through its hidden state update (equation \ref{eq:vanilla-rnn-equation}).
At a given timestamp $t = \overline{1, T}$, for an input $x_t$, the hidden state $h_t$ is computed as: 
\begin{equation}
h_t = \phi(W_{xh}x_t + W_{hh}h_{t-1} + b_h),
\label{eq:vanilla-rnn-equation}
\end{equation}

where $W_{xh}$ and $W_{hh}$ are learnable weight matrices, $b_h$ is a bias term, and $\phi(\cdot)$ 
denotes a non-linear activation function such as $\tanh$ or sigmoid. The hidden state is then used to 
produce the output $y_t$: 

\begin{equation}
y_t = W_{hy}h_t + b_y,
\label{eq:y_t RNN}
\end{equation}

where $W_{hy}$ and $b_y$ are the output weight matrix and bias, respectively. Through the recurrent 
connection involving $h_{t-1}$, the network is able to propagate information across time steps, 
allowing previous inputs to influence future predictions.\\ \\ 
%
Training recurrent neural networks requires propagating gradients not only through the layers of the network but also 
through time. This is typically achieved using \textit{Backpropagation Through Time} (BPTT), 
introduced by Werbos et al.\cite{werbos1990backpropagation}. BPTT unfolds the recurrent network across the input sequence, 
transforming it into an equivalent deep feedforward network with shared parameters at each time step. Gradients
are then computed by applying the backpropagation algorithm to the unfolded computational graph and accumulated 
across all time steps before updating the model parameters. \\ \\ 
%
While this training procedure allows RNNs to learn temporal dependencies, it also exposes a major limitation of the 
architecture. During BPTT, gradients are propagated backward through all previous time steps of the sequence. 
At each step, the gradient is repeatedly multiplied by the same recurrent transformation, as it is passed through the 
unfolded network in time. As a result, the gradient magnitude may gradually decrease or increase with each multiplication.
When it becomes too small, the network struggles to learn from earlier time steps, a phenomenon known as the \textit{vanishing gradient problem}. 
Conversely, when it grows too large, training becomes unstable, leading to the \textit{exploding gradient problem} \cite{bengio1994learning}.\\ \\ 
%
These issues significantly limit the ability of standard RNNs to capture long-range dependencies in sequential data. 
To address the exploding gradient problem, Pascanu et al. \cite{pascanu2013difficulty} introduced \textit{gradient clipping}, 
a technique that limits the magnitude of gradients during training and improves optimization stability. Nevertheless, 
the difficulty of preserving information over long sequences remained a fundamental challenge, motivating the development 
of more advanced recurrent architectures such as Long Short-Term Memory (LSTM)\cite{10.5555/2998981.2999048} networks 
and Gated Recurrent Units (GRUs)\cite{cho2014learningphraserepresentationsusing}\cite{chung2014empiricalevaluationgatedrecurrent}.

\subsection{Long-Short Term Memory}

The \textbf{Long Short-Term Memory (LSTM)} architecture, introduced by Hochreiter and Schmidhuber \cite{10.5555/2998981.2999048} in the 1990s,
was proposed to address the limitations of standard RNNs when learning long-range dependencies. The key innovation of the architecture is the 
introduction of a dedicated memory state that can preserve information across many time steps. Instead of relying solely on the hidden state, 
the LSTM uses a set of gating mechanisms to control how information is written to, stored in, and retrieved from memory. 
This allows the model to memorize relevant information over longer sequences while reducing the impact of the vanishing gradient problem 
encountered in traditional recurrent networks. \\ \\
%
\begin{figure}[h!]
    \centering
    \includegraphics[width=\textwidth]{LSTM.png}
    \caption{The LSTM cell, figure inspired by \cite{olah2015lstm}}
    \label{fig:LSTM}
\end{figure}

\noindent The standard LSTM cell maintains an internal cell state, $c_t$(\ref{eq:lstm_cell_state}), which acts as a long-term memory buffer, 
and a hidden state, $h_t$(\ref{eq:lstm_hidden_state}), 
which serves as the short-term working memory routed to subsequent layers. The flow of information inside the cell is controlled by 
three gates with different scopes. The forget gate(\ref{eq:lstm_forget_gate}) manages what proportion of the historical cell state to discard by analyzing the 
current input and the previous hidden state. The input gate(\ref{eq:lstm_input_gate}) controls how much new information, computed from the current input and 
previous hidden state, is incorporated into the cell state. Finally, the output gate(\ref{eq:lstm_output_gate}) determines which information stored in the 
updated cell state is propagated as the next hidden state.

\begin{equation}
f_t = \sigma(W_f x_t + R_f h_{t-1} + b_f)
\label{eq:lstm_forget_gate}
\end{equation}

\begin{equation}
i_t = \sigma(W_i x_t + R_i h_{t-1} + b_i)
\label{eq:lstm_input_gate}
\end{equation}

\begin{equation}
\tilde{c}_t = \tanh(W_c x_t + R_c h{t-1} + b_c)
\label{eq:lstm_candidate_state}
\end{equation}

\begin{equation}
c_t = f_t \odot c_{t-1} + i_t \odot \tilde{c}_t
\label{eq:lstm_cell_state}
\end{equation}

\begin{equation}
o_t = \sigma(W_o x_t + R_o h_{t-1} + b_o)
\label{eq:lstm_output_gate}
\end{equation}

\begin{equation}
h_t = o_t \odot \tanh(c_t)
\label{eq:lstm_hidden_state}
\end{equation}

\noindent Overall, the introduction of a dedicated memory cell and gating mechanisms 
allows LSTMs to address the limitations of standard RNNs in capturing long-range dependencies. 
However, the sequential processing over the time dimension still limits parallelization and makes them 
computationally expensive for long sequences. Several architectural extensions were later 
proposed to improve their effectiveness in different settings. Bidirectional recurrent networks \cite{schuster1997bidirectional} 
improved contextual understanding by processing sequences in both forward and backward directions,
while encoder-decoder architectures \cite{sutskever2014sequence} enabled sequence-to-sequence learning by separating the encoding 
and decoding processes.
In addition, stacked recurrent architectures \cite{graves2013generating} increased model capacity 
by composing multiple recurrent layers, allowing for more expressive hierarchical representations.\\ \\ 
%
Despite these improvements, these variants did not remove the fundamental sequential bottleneck of recurrent 
computation, which limited scalability on long sequences. These limitations motivated the transition toward 
architectures that avoid recurrence entirely, most notably Transformer\cite{vaswani2017attention} models based on self-attention. 
In this broader context, recent work has revisited recurrent memory mechanisms at scale, leading to 
architectures such as xLSTM, which aim to combine the efficiency of structured recurrence with improved 
scalability and expressive memory modeling. \\ \\ 

\subsection{Extended Long-Short Term Memory}

The extended LSTM (xLSTM) architecture was introduced by Beck et al. \cite{beck2024xlstmextendedlongshortterm} 
to overcome these bottlenecks by leveraging modern techniques 
while fundamentally modifying the LSTM memory structure. 

\subsubsection{Addressing Classic LSTM Limitations}
To understand the architectural shifts in the xLSTM, it is necessary to outline the 
three main bottlenecks of the classic LSTM:
\begin{enumerate}
    \item \textbf{Inability to revise storage decisions:} Traditional LSTMs struggle to
    update or revise a stored value when new, more relevant information is found. 
    %
    \item \textbf{Limited storage capacities:} Information in a classic LSTM is compressed 
    into a scalar cell state, leading to poor performance when predicting rare tokens or 
    retaining extensive context.
    %
    \item \textbf{Lack of parallelizability:} The recurrent memory update relies on hidden-to-hidden 
    connections between consecutive time steps, requiring strictly sequential processing and 
    preventing the highly parallel computation that characterizes Transformer architectures.
\end{enumerate}

\noindent To address these issues, xLSTM introduces several changes to the classical LSTM architecture.
First, it replaces the traditional logistic gating mechanism with \textit{exponential gating}, 
which provides a less restricted range for controlling memory updates. To ensure numerical stability 
during training and inference, xLSTM uses normalization and 
stabilization techniques. In addition, the architecture looks to improve the design of the memory state, 
resulting in two distinct recurrent building blocks: the \textit{scalar LSTM (sLSTM)} and the 
\textit{matrix LSTM (mLSTM)}. While sLSTM extends the traditional LSTM memory formulation with 
improved gating and memory management, mLSTM introduces a matrix-based memory structure capable of 
storing richer associations between tokens.

\subsubsection{sLSTM}
The sLSTM (scalar LSTM) remains structurally closer to the classic architecture but 
introduces exponential activation functions for the input and forget gates. 
This modification allows the model to control its storage more effectively:

% Ensure you have \usepackage{amsmath} and \usepackage{bm} in your preamble
\begin{align}
    \mathbf{c}_t &= \mathbf{f}_t \odot \mathbf{c}_{t-1} + \mathbf{i}_t \odot \mathbf{z}_t && && \text{cell state} \\
    \mathbf{h}_t &= \mathbf{o}_t \odot \tilde{\mathbf{h}}_t, & \tilde{\mathbf{h}}_t &= \psi(\mathbf{c}_t) && \text{hidden state} \\
    \mathbf{z}_t &= \varphi(\tilde{\mathbf{z}}_t), & \tilde{\mathbf{z}}_t &= \boldsymbol{W}_z \boldsymbol{x}_t + \boldsymbol{R}_z \boldsymbol{h}_{t-1} + \boldsymbol{b}_z && \text{cell input} \\
    \mathbf{i}_t &= \sigma(\tilde{\mathbf{i}}_t), & \tilde{\mathbf{i}}_t &= \boldsymbol{W}_{\mathbf{i}} \boldsymbol{x}_t + \boldsymbol{R}_{\mathbf{i}} \boldsymbol{h}_{t-1} + \boldsymbol{b}_{\mathbf{i}} && \text{input gate} \\
    \mathbf{f}_t &= \sigma(\tilde{\mathbf{f}}_t), & \tilde{\mathbf{f}}_t &= \boldsymbol{W}_{\mathbf{f}} \boldsymbol{x}_t + \boldsymbol{R}_{\mathbf{f}} \boldsymbol{h}_{t-1} + \boldsymbol{b}_{\mathbf{f}} && \text{forget gate} \\
    \mathbf{o}_t &= \sigma(\tilde{\mathbf{o}}_t), & \tilde{\mathbf{o}}_t &= \boldsymbol{W}_{\mathbf{o}} \boldsymbol{x}_t + \boldsymbol{R}_{\mathbf{o}} \boldsymbol{h}_{t-1} + \boldsymbol{b}_{\mathbf{o}} && \text{output gate}
\end{align}

\noindent To prevent the numerical overflows specific to the exponential function, a normalizer 
state ($n_t$) and stabilization states are utilised:
\begin{align}
    m_t &= \max \left( \log(f_t) + m_{t-1}, \log(i_t) \right) && \text{stabilizer state} \label{eq:stabilizer-state}\\
    i'_t &= \exp \left( \log(i_t) - m_t \right) = \exp(\tilde{\imath}_t - m_t) && \text{stabil. input gate} \\
    f'_t &= \exp \left( \log(f_t) + m_{t-1} - m_t \right) && \text{stabil. forget gate}
\end{align}

\noindent Despite these changes, because it retains hidden-hidden recurrent connections, 
the sLSTM doesn't improve the sequential processing bottleneck of the classic LSTM.
For the scope of this work, the architecture used in the experiments solely relies on mLSTM cells. 

\subsubsection{mLSTM}
The mLSTM fundamentally shifts the LSTM paradigm to solve both the storage capacity 
and parallelization bottlenecks. To overcome the limited capacity of scalar states, the mLSTM 
expands the memory cell into a matrix $C \in \mathbb{R}^{d \times d}$ where $d$ represents the embedding dimension of $x_{t}$. Crucially, to achieve 
this efficiency, the mLSTM entirely abandons memory mixing (the hidden-hidden recurrent connections) allowing parallelization. \\ \\ 
%
Relying on Bidirectional Associative Memory (BAM)\cite{kosko1987adaptive} principles and Transformer terminology,
the mLSTM stores information as key-value pairs ($k_t$, $v_t$) using a covariance (outer product) 
update rule, and retrieves it using a query vector ($q_t$). The mLSTM cell also uses the stabilizer state introduced in the sLSTM equations(\ref{eq:stabilizer-state})
%
\begin{align}
    \boldsymbol{C}_t &= \mathrm{f}_t \boldsymbol{C}_{t-1} + \mathrm{i}_t \boldsymbol{v}_t \boldsymbol{k}_t^\top && && \text{cell state} \\
    \boldsymbol{n}_t &= \mathrm{f}_t \boldsymbol{n}_{t-1} + \mathrm{i}_t \boldsymbol{k}_t && && \text{norm. state} \\
    \boldsymbol{h}_t &= \mathbf{o}_t \odot \tilde{\boldsymbol{h}}_t, & \tilde{\boldsymbol{h}}_t &= \boldsymbol{C}_t \boldsymbol{q}_t / \max \left\{ \left| \boldsymbol{n}_t^\top \boldsymbol{q}_t \right|, 1 \right\} && \text{hidden state} \\
    \boldsymbol{q}_t &= \boldsymbol{W}_q \boldsymbol{x}_t + \boldsymbol{b}_q && && \text{query input} \\
    \boldsymbol{k}_t &= \frac{1}{\sqrt{d}} \boldsymbol{W}_k \boldsymbol{x}_t + \boldsymbol{b}_k && && \text{key input} \\
    \boldsymbol{v}_t &= \boldsymbol{W}_v \boldsymbol{x}_t + \boldsymbol{b}_v && && \text{value input} \\
    \mathrm{i'}_t &= \exp(\tilde{i}_t - m_t), & \tilde{i}_t &= \boldsymbol{w}_{\mathrm{i}}^\top \boldsymbol{x}_t + b_{\mathrm{i}} && \text{input gate} \\
    \mathrm{f'}_t &= \exp \left( \log(\sigma(\tilde{f_t})) + m_{t-1} - m_t \right), & \tilde{f}_t &= \boldsymbol{w}_{\mathrm{f}}^\top \boldsymbol{x}_t + b_{\mathrm{f}} && \text{forget gate} \\
    \mathbf{o}_t &= \sigma(\tilde{\mathbf{o}}_t), & \tilde{\mathbf{o}}_t &= \boldsymbol{W}_{\mathbf{o}} \boldsymbol{x}_t + \boldsymbol{b}_{\mathbf{o}} && \text{output gate}
\end{align}

\subsection{The xLSTM Block}

The xLSTM block leveages the mLSTM as a multi-head sequence-mixing layer. First, the input token 
representation $x_t$ is normalized with RMSNorm. The normalized vector is then linearly projected into several heads.
Each head receives its own projected query, key, value, and gate information, and each head runs an independent mLSTM cell. 
Therefore, the recurrent state is also stored separately for each head. The outputs of the heads, written as $h_1$, $h_2$, ..., 
$h_N$, are then normalized per head, gated, concatenated, and passed through a shared output linear layer. This produces a vector 
with the same hidden dimension as the original input, so it can be added back through a residual connection. \\ \\ 
%
After the mLSTM sub-layer, the block applies a second pre-RMSNorm residual sub-layer: a SwiGLU feed-forward network. 
This layer does not mix information across time. Instead, it processes each token representation independently. 
The final output of the block is obtained by adding the FFN output back to the result of the mLSTM residual branch. \\ \\
%
\begin{figure}[h!]
    \centering
    \includegraphics[width=\textwidth]{mlstm-block.drawio.png}
    \caption{mLSTM Block}
    \label{fig:mLSTM-block}
\end{figure}
\noindent This block structure is the one used in the available pre-trained xLSTM-7b model. It differs from the original 
xLSTM paper, where the mLSTM block used a more specialized structure with an up-projection, causal convolution,
block-diagonal query/key projections, and an internal skip connection. In the pre-trained model, the block is 
simpler and closer to a standard transformer block: a multi-head mLSTM sequence-mixing layer followed by a 
SwiGLU feed-forward layer. For the scope of this study, we'll refer to this version of the xLSTM layer.